% THETA3-B19 contribution by /root/definite_bridge, 21 November 2026.
% Static source reading; no Lean build is asserted.
\subsection{THETA3-B19: definite periods are finite CM theta-jet functionals}\label{B19}

The remaining arithmetic comparison can be localized to explicit CM
abelian varieties.  This gives a common period normalization for a
fixed finite family of inputs and source characters.  It does not
assert that this normalization equals the particular constant of an
uninstantiated \texttt{MainX} datum.

\paragraph{The actual CM abelian varieties at the evaluation point.}
Keep the rank-three setting of 032 and B12.  Choose a rational
negative line $L=Ez$ and an orthogonal decomposition
\[
 V=Ee_1\perp Ee_2\perp Ez
\]
over $E$, with the first two lines positive at the active place.
They are positive at the other real places as well.  For the
multiplicity line $P_i$, put
\[
 d_i(\tau)=v_i(\tau^{-1}),\qquad
 d_i^{\mathrm{flip}}(\tau)=
 \begin{cases}-d_i(\tau)&\tau\in\{\mu,c\mu\},\\
 d_i(\tau)&\text{otherwise}.
 \end{cases}
\]
Use the geometric $E$-action on
$A_{x,i}=\mathcal F_i^\vee/\iota(\Lambda_i)$, where
$\iota(u)(v)=\psi_i(v,u)$ and $x=[Ez]$.
Then, as rational $E$-Hodge structures,
\begin{equation}\label{B19-CM}
 H^1(A_{x,i},\Q)
 \simeq U(d_i)^{\oplus2}\oplus U(d_i^{\mathrm{flip}}),
 \qquad
 A_{x,i}\sim_E B(d_i)^2\times B(d_i^{\mathrm{flip}}).
 \tag{B19a}
\end{equation}
Here the last relation is an isogeny, with the chosen polarizations
retained separately.  In particular it is not an assertion that the
polarization is the product principal polarization.

To prove the formula, use 032, \S4.2, equation (4.5), pp.14--15.
For the positive coordinate lines the homology-filtration block
uses $c\tau$ when $d_i(\tau)=1$ and $\tau$ when
$d_i(\tau)=-1$, including the active place.  The negative
coordinate agrees at every compact place and switches the active
pair.  The relation
\[
 \iota(au)(v)=\iota(u)(\overline a v)
\]
conjugates these embedding labels in geometric $H^1$.
This proves the first formula.  The equivalence between rational
weight-one Hodge maps and homomorphisms of abelian varieties up to
isogeny, used in 032, Lemma 2.2, proves the second.
The two-dimensional tangent block used by the first Taylor
coefficient comes from the two $B(d_i)$ factors.  The block has
the archive tangent label $\mu$, whereas the covariant Taylor
tensor has label $c\mu$, as in A8.  The inverse in the definition
of $d_i$ cannot in general be discarded.

The descent input is Milne, \emph{Complex Multiplication},
version 14 July 2020, Proposition 7.9 and Corollary 7.10,
printed pp.54--55: extension from an algebraically closed subfield
of $\C$ is fully faithful for abelian varieties and is essentially
surjective on CM abelian varieties.  Apply it to
$\overline\Q\subset\C$.  The abelian variety in (B19a), together with its polarization,
symmetric theta line, and any specified finite torsion data, has a
model over $\overline\Q$.  This assertion needs no identification
of its polarization with a product polarization.  Start with the
CM factors $B(d)$ over $\overline\Q$.  A rational Hodge
isomorphism gives a complex isogeny; its kernel consists of torsion
points and is defined over a finite extension of their field of
definition.  Taking the quotient produces the required abelian
variety over $\overline\Q$.  Homomorphisms and polarizations of
abelian varieties over an algebraically closed field of
characteristic zero are unchanged by extension to $\C$; thus the
given polarization descends as well.  Its symmetric line-bundle
representatives differ by two-torsion data, which are also
algebraic.  These operations preserve the specified analytic
polarized torus and merely supply its algebraic model.

\paragraph{The finite theta module and its one common comparison scalar.}
Choose rational symplectic coordinates for the actual alternating
form.  Replacing a lattice by a commensurable symplectically
self-dual lattice is an isogeny and allows principal coordinates
without changing the rational CM decomposition.  Record this
isogeny and its effect on the polarization; it is part of the
comparison data.  Write the resulting period matrix as $Z_i$ and
the torus as $\C^g/(\Z^g+Z_i\Z^g)$, where $g=3[F:\Q]$.

For a fixed finite family of finite Schwartz inputs with algebraic
values, including
all the finite translates used below, choose an even integer $N_i$
such that their supports are contained in $N_i^{-1}\Z^g$ and
their restrictions to rational points are invariant under
$N_i\Z^g$.  Enlarge $N_i$ to absorb the chosen isogenies.  They
belong to the finite Schr\"odinger space
\[
 W_i=\operatorname{Fun}
       (N_i^{-1}\Z^g/N_i\Z^g,\overline\Q).
\]
The theta sum followed by $z=N_iw$ gives the section
\[
 \mathcal T_i(\varphi)(w)=
 \exp\!\left(\frac\pi2N_i^2w^t
                       (\operatorname{Im}Z_i)^{-1}w\right)
 \sum_{q\in N_i^{-1}\Z^g}
    \varphi(q)\exp(\pi i q^tZ_iq+2\pi iN_iq^tw)
 \tag{B19b}
\]
of the $N_i^2$-th power $L_i^{N_i^2}$ of the symmetric theta
line.  In restriction-of-scalars coordinates the exponents are
the corresponding sums over embeddings.  Any scalar in the
chosen oscillator or Fourier normalization is included in the
definition of $\mathcal T_i$, not suppressed.

The dimension of $W_i$ is $N_i^{2g}$, equal to
$h^0(A_{x,i},L_i^{N_i^2})$.  The map $\mathcal T_i$ is injective:
its Fourier expansion recovers every value of $\varphi$ on the
finite quotient.  It is therefore an isomorphism after extension
to $\C$.  Translation and multiplication by finite characters
give its theta-group equivariance, with the central character
fixed by the displayed multiplier.  The completion factor is
precisely the one proved in A8.

Choose the matching algebraic theta structure and an equivariant
isomorphism over $\overline\Q$,
\[
 J_i:W_i\longrightarrow
       H^0(A_{x,i},L_i^{N_i^2}).
\]
The matching theta structure is algebraic: its torsion points are
algebraic, and lifts of finite order in the theta group satisfy
algebraic power equations.  A choice of compatible finite lifts
and of the central roots of unity therefore descends after a
finite extension.  It is essential to match this theta structure,
not merely the polarization type.

There is a single $c_i\in\C^\times$ such that
\begin{equation}\label{B19-Heis}
                 \mathcal T_i=c_i(J_i\otimes\C).
                 \tag{B19c}
\end{equation}
For completeness, the finite Heisenberg proof is elementary.
Multiplication characters diagonalize $W_i$ with one-dimensional
eigenspaces spanned by the delta functions.  Translation operators
permute these eigenspaces transitively.  A commuting endomorphism
is diagonal in the delta basis, and commutation with translations
forces all its diagonal entries to agree.  The same representation
is the theta-group representation on the section space.  This is
Mumford, \emph{On the equations defining abelian varieties I},
\S1, Proposition 3 and Theorem 2, printed pp.295--298.
Applying the commuting-endomorphism argument to
$\mathcal T_iJ_i^{-1}$ proves (B19c).

Thus $c_i$ depends on the actual torus, theta structure, oscillator
normalization, and chosen algebraic isomorphism $J_i$.  It does
not depend on $\varphi$.  It may be specified, without an
unmentioned scalar, by choosing a section with nonzero value at
the origin and taking the ratio of its two values in the fixed
rigidification.  Such a section exists because $N_i^2\geq4$.
This ratio is a theta-null period; no assertion that $c_i$ itself
is algebraic is made.

\paragraph{First jets and the CM tangent periods.}
Normalize the symmetric structure on $L_i^{N_i^2}$ to act as the
identity on its fibre at zero, and fix an algebraic rigidification.
A symmetric formal local generator with value one exists by
averaging a generator with its pullback by $[-1]$.  The first jet
of a section in this generator is independent of the choice:
the ratio of two such generators is an even unit with value one,
so its differential at zero vanishes.  First jets of algebraic
sections are therefore algebraic in algebraic tangent frames.
This argument concerns first jets.  A8 supplies the specified
canonical gauge for higher jets.

The Appell--Humbert frame in (B19b) has the same normalized
symmetry.  Its comparison with the symmetric algebraic frame has
value one and zero first derivative at the origin.  Consequently
(B19c) gives, for every input, the exact first-jet comparison.
Let $u^z_{ir}$, $r=1,2$, be the two directions in the original
$z$ variable used by the selected $c\mu$ Taylor tensor.  Let
$u^w_{ir}$ have the same coordinate vectors in the $w$ variable;
thus $d(z=N_iw)(u^w_{ir})=N_i u^z_{ir}$.  It is the
$w$ variable that uniformizes the chosen algebraic torus.  Choose
algebraic tangent vectors $w_{i1},w_{i2}$ through the two CM
factors in (B19a).  There exist $p_i\ne0$ and
$M_i\in\operatorname{GL}_2(\overline\Q)$ such that
\[
 u^w_{ir}=p_i\sum_{s=1}^2 (M_i)_{sr}w_{is}.
 \tag{B19d}
\]
Here $p_i$ is defined by this Betti--de Rham tangent comparison,
so the equation fixes whether a period or its inverse occurs.
To justify the common scalar, use the same algebraic frame on
the one-dimensional $\mu$ tangent eigenspace of $B(d_i)$ in
each of its two copies.  The two copies then have the same
comparison factor.  The rational $E$-linear isogeny and the
chosen rational coordinate lines change their frames by matrices
over $\overline\Q$, giving $M_i$.  This proves (B19d).

If $j_{is}(\varphi)$ denotes the algebraic first jet of
$J_i\varphi$ along $w_{is}$, the result is
\begin{equation}\label{B19-jet}
 D_{u^z_{ir}}\widehat\vartheta_{\varphi}(Z_i,0)
   =\frac{c_ip_i}{N_i}
        \sum_{s=1}^2(M_i)_{sr}j_{is}(\varphi).
 \tag{B19e}
\end{equation}
Indeed $D_{u^w_{ir}}\mathcal T_i(\varphi)(0)
=N_iD_{u^z_{ir}}\widehat\vartheta_{\varphi}(Z_i,0)$.
The factor $1/N_i$ therefore occurs once.  At first order the
quadratic completion contributes zero, and the divided Taylor
coefficient has $1/1!=1$.  If a different rational symplectic
lattice is used, its pullback matrix and its rational dilation
must be included in $M_i$ and $N_i$.

The passage from the two selected Taylor tensors to
$\Omega_x^1$ uses the chosen square-root determinant and
constant-line frames of 032, equation (4.6).  In a fixed rational
coordinate frame of $V$ its remaining scalar for the wedge will
be denoted by $\kappa_x\ne0$.  This scalar is part of the
specified metaplectic and canonical-fibre comparison.  All
algebraic coordinate matrices can be absorbed into $M_1,M_2$.
Define the common CM comparison factor
\[
 C_x=\kappa_x\frac{c_1c_2p_1p_2}{N_1N_2}.
 \tag{B19f}
\]
No Haar volume is included in $C_x$.

\paragraph{The infinity integral and the finite source class sums.}
Let $T_i=\U(P_i)$.  Its entire archimedean group is compact.
Fix a finite open subgroup $K_i$ fixing all source inputs, and
characters $\pi_i$ trivial on $K_i$.  The chosen first-degree
archimedean oscillator coefficient has one character
$\eta_{i,\infty}$ under $T_i(F_\infty)$.  Thus the bilinear
integral of the theta coefficient against $\pi_i$ vanishes unless
\[
                    \eta_{i,\infty}\pi_{i,\infty}=1.
 \tag{B19g}
\]
Indeed the Haar integral of a nontrivial character of a compact
group is zero.  Assume (B19g), as is necessary for a nonzero
degree-one lift in this fixed infinity type.  The integrand is
then invariant under the full archimedean group and under $K_i$.
Choose representatives $t_{ia}$ of the finite quotient
$T_i(F)\backslash T_i(\A_{F,f})/K_i$ and put
\[
 \Gamma_{ia}=T_i(F)\cap t_{ia}K_it_{ia}^{-1},\qquad
 v_i^T=\operatorname{vol}(T_i(F_\infty))\operatorname{vol}(K_i).
\]
The intersection defining $\Gamma_{ia}$ is understood with the
archimedean component unrestricted; it is a finite group of
integral norm-one units.  Quotient Haar integration gives exactly
\[
 \int_{[T_i]}\theta_i(g,t)\pi_i(t)\,dt
  =v_i^T\sum_a\frac{\pi_{i,f}(t_{ia})}{|\Gamma_{ia}|}
                      \theta_i(g,t_{ia}).
 \tag{B19h}
\]
This is a consequence of the verified archimedean character
cancellation, not a replacement of an arbitrary continuous
integral by a finite sum.  Stabilizers and both Haar factors are
retained.  For an unintegrated Taylor form set $v_i^T=1$ and
use its single prescribed input instead of (B19h).

For honest algebraic infinity characters the finite values in
(B19h) are algebraic.  Here is a direct proof.  The finite torus
class group has finite order.  For a representative $t$ some
$r>0$ satisfies $t^r=\gamma k$ with
$\gamma\in T_i(F)$ and $k\in K_i$.  Global triviality on
$T_i(F)$ gives
\[
 \pi_{i,f}(t)^r=\pi_{i,\infty}(\gamma_\infty)^{-1}
                      \in\overline\Q^\times.
\]
The last inclusion follows because an infinity character of the
honest norm-one torus is a product of integral powers of its
algebraic embedding coordinates.  Taking a root proves the claim.
Weil transforms of the rational finite inputs are likewise
$\overline\Q$-valued when the fixed finite splitting characters
are algebraic: translations and character multiplication use
algebraic values, and a finite Fourier transform uses roots of
unity and the specified algebraic self-dual volume factors.
Arbitrary scalar multiples of the input must be carried as
additional coefficients, rather than declared algebraic.

\paragraph{The definite Bessel coefficient is the promised finite functional.}
Let $H=\U(L^\perp)$ and choose a common finite level $K_H$.
Let $e$ be a function in its scalar archimedean type, the same
type as the canonical coefficient $f$ of the wedge.  Choose
representatives $h_b$ of
$H(F)\backslash H(\A_{F,f})/K_H$ and set
\[
 \Gamma_b^H=H(F)\cap h_bK_Hh_b^{-1},\qquad
 v_H=\operatorname{vol}(H(F_\infty))\operatorname{vol}(K_H).
\]
Use these fixed Haar measures in B17 as well.  The product
$f\overline e$ is invariant under $H(F_\infty)K_H$, so
\[
 \mathcal B_e(f)=\int_{[H]}f(h)\overline{e(h)}\,dh
       =v_H\sum_b\frac{f(h_b)\overline{e(h_b)}}{|\Gamma_b^H|}.
 \tag{B19i}
\]
At $h_b$ the finite oscillator input in the $a$-summand is the
actual transform $\varphi_{iab}=\omega_{i,f}(h_b,t_{ia})
\varphi_i$.  The common level $N_i$ above was chosen to include
all these finitely many transforms.  Define the algebraic arrays
\[
 a_{ibr}=\sum_a\frac{\pi_{i,f}(t_{ia})}{|\Gamma_{ia}|}
             \sum_{s=1}^2(M_i)_{sr}j_{is}(\varphi_{iab}).
\]
Equations (B19e)--(B19h), and alternation of the two one-forms,
give the exact identity
\begin{equation}\label{B19-period}
 \mathcal B_e(f)
   =v_Hv_1^Tv_2^T C_x
       \sum_b\frac{\overline{e(h_b)}}{|\Gamma_b^H|}
          (a_{1b1}a_{2b2}-a_{1b2}a_{2b1}).
 \tag{B19j}
\end{equation}
The sign is the order of the chosen coordinate volume form and
of the two input one-forms.  No factorial is missing: this is the
alternation of two first Taylor tensors.

Choose $e$ from the algebraic finite definite module, meaning that
its values in the specified algebraic scalar-type frame belong
to $\overline\Q$.  Then the sum in (B19j) is algebraic.  For
a fixed finite family of inputs, characters and test vectors,
it lies in one number field.  The same $C_x$, $v_H$ and $v_i^T$
work throughout that family, provided the types, theta structures,
levels, splittings, and frames remain fixed.  This proves the
absence of an extra transcendental scalar depending on the
source character in this finite CM-evaluation comparison.

An orthonormal definite basis need not have the same algebraic
normalization.  Its Gram matrix in an algebraic basis is $v_H$
times a matrix with algebraic entries, because the stabilizer
weights are rational.  Gram--Schmidt therefore gives basis values
in $v_H^{-1/2}\overline\Q$.  For that normalization the common
measure factor in (B19j) is $v_H^{1/2}v_1^Tv_2^T C_x$.
This records the change instead of silently imposing Haar
probability or identifying an algebraic vector with a unit vector.

There is also an explicit lower certificate once arithmetic data
for the displayed sum are supplied.  Let its value be $b\ne0$
in a number field $K\subset\C$ of degree $d$ whose chosen
embedding is nonreal (adjoin $i$ if necessary), let $Db$ be integral
for a positive integer $D$, and assume
$|\sigma(b)|\leq M$ with $M>0$ for the $d-2$ embeddings other than the
chosen complex embedding and its conjugate.  Then
\[
 |b|^2\geq D^{-d}M^{-(d-2)},\qquad
 |\mathcal B_e(f)|^2\geq
 |v_Hv_1^Tv_2^TC_x|^2D^{-d}M^{-(d-2)}.
 \tag{B19k}
\]
Indeed $|N_{K/\Q}(Db)|\geq1$, and the two distinguished
factors in the norm product are $D^2|b|^2$.
For a unit test vector replace the Haar factor as just explained.
This certificate does not supply a bound on the CM period factor
or the conjugates without additional data.

\paragraph{What remains to match the actual archive arithmetic functional.}
The recovered source makes this remaining requirement explicit.
Under the directory
\texttt{Development/AbelianVarietyTheory/Definitions/LatticeFrames/},
\texttt{LevChiEqOne.lean:73--77} defines
\begin{verbatim}
def abelian₂ ... :=
  ((I.chIn u a).rho i c * (I.chIn u a).rho' i c') •
    jetL I.J.cD
      (I.th u a i _ _ ((I.stp u a i).rep c)
                         ((I.stp u a (i + 1)).rep c'))
\end{verbatim}
The omitted arguments are only the typed indices and the declared
domain; the scalar and section in this quotation are the actual
definition.  Thus this functional is tied to a specified
arithmetic theta section and its period frame.
\texttt{OnFramesRdApply.lean:108--113} requires its equality
with the analytic theta kernel after multiplication by
\texttt{volLev}, the prescribed kernel constant \texttt{t.cK},
and the polynomial rescaling \texttt{rho.sc}.

The proved constructor is
\texttt{Development/AbelianVarietyTheory/Proofs/DecoratedLatticeProofs/KernelPartLawsOfColumnPins.lean},
theorem \texttt{clsKL\_is}, lines 22--52.  Its core step at
lines 29--35 invokes
\texttt{thetaKernelIdentity\_at\_every\_tuple} with
\texttt{PinX} and \texttt{JetS}.
The underlying proved reading is
\texttt{ThetaKernelSplitPlacesEqJetClassicalTheta.lean},
theorem \texttt{thetaKernel\_eq\_jet\_thC}, lines 46--71.
Its \texttt{StdPair} hypothesis is defined in
\texttt{SplitPart.lean:38--49}; it includes
\begin{verbatim}
gS : S.gram = 1
gT : T.gram = 1
rk : 0 < S.rank ∧ S.rank < T.rank ∧ Even S.rank ∧ Even T.rank
cc : D.host.ccK = ⊤
\end{verbatim}
The explicit jet scale is exhibited in
\texttt{PureTensorOZeroEqJetP.lean:342--346} as
$s=c_K(-\delta)$, where lines 267--269 define
$c_K=(\sqrt{2\pi/\sqrt d})^{-1}$.
These are static source readings at the pinned mirror commit
recorded in B0, not a build or an audit of their full dependencies.

Consequently the arithmetic match for (B19j) is now a concrete
comparison of the polarized CM lattices and isogenies, finite
theta structures and sections, the two tangent period frames,
the determinant-line frame, and the displayed covolume/Haar
factors with the corresponding \texttt{I.th}, \texttt{I.J.cD},
\texttt{rho}, \texttt{c0}, and \texttt{t.cK}.  It is sufficient
to check equality on a basis of the finite theta modules and the
relevant first-jet directions; linearity then gives the entire
fixed family and (B19j).  The recovered definite Gram-one,
strictly increasing even-rank instance does not itself make that
comparison for the rank-two/rank-three mixed-signature data.
Defining an abstract \texttt{abelian₂} to equal the analytic side
would satisfy a formal equality but would not establish this
arithmetic comparison.

In particular (B19j) gives an actual definite Bessel functional
of the CM jets associated with the original 032 inputs, and B17
identifies the same analytic Bessel coefficient with the last
theta seesaw.  The remaining scalar comparison is common to
the fixed family, with every measure and frame identified; it is
not a collection of unspecified constants indexed by characters.
\hfill$\square$
